\chapter{Modello ad un fluido (Magnetoidrodinamica)}

Assumendo che si trovi una chiusura, come l'incomprimibilità o una politropica, il nostro sistema di equazioni fluide a cui ci siamo ricondotti nel capitolo precedente vanno ancora accoppiate alle equazioni dell'elettrodinamica.

Finora abbiamo trovato le equazioni di bilancio del numero di particelle (equazione di continuità), equazione della quantità di moto (o impulso) ed equazione dell'energia per le varie specie che compongono il plasma, ma nel caso più semplice di due fluidi (ioni + elettroni) possiamo fare alcune definizioni per ricondurci ad una descrizione a un fluido, la \textbf{magnetoidrodinamica}; descrizione che spazia su molti ordini di grandezza ma perdendo alcuni aspetti cinetici. Questa teoria è valida in condizioni come 
\begin{itemize}
	\item il moto è fortemente collettivo $\iff$ vi devono essere molte particelle in un volume di Debye ($n\, \lambda_D^2 \gg 1$);
	\item collisioni abbastana frequenti da mantenere il sistema in equilibrio locale;
	\item lunghezze caratteristiche $\gg \lambda_D$;
	\item frequenze caratteristiche $\ll \omega_P$.
\end{itemize}

Diamo come prima cosa alcune definizioni per il singolo fluido: 

\[ 
\begin{split}
	\textbf{densità di massa} &:\ \ \ \ \ \ \ \ \ \ \ \ \rho\ =\ \sum_\alpha m_\alpha n_\alpha = m_i n_i + m_e n_e \sim \frac{n}{2}(m_i + m_e) \\
	\textbf{velocità media} &: \ \ \ \ \ \ \ \ \ \ \ \ \bar u\ =\ \frac{\sum_\alpha m_\alpha n_\alpha u_\alpha}{\sum_\alpha m_\alpha n_\alpha}\sim \frac{n_0}{n_0} \frac{m_i \vec u_i + m_e \vec u_e }{m_i + m_e} \sim \vec u_i + \frac{m_e}{m_i} \vec u_e\\
	\textbf{corrente} &:\ \ \ \ \ \ \ \ \ \ \ \ \vec J\ =\ \sum_\alpha q_\alpha n_\alpha \vec u_\alpha
\end{split}
\]

Si nota che per la densità di massa il contributo è dato principalmente dagli ioni, come per la velocità dove al primo ordine potremmo approssimare $\bar u \sim \vec u_i$; tuttavia, costruire la teoria a un fluido non è sempre semplice come in quest'ultima approssimazione. Andiamo a vedere innanzitutto come cambiano le nostre equazioni fluide se sommate assieme ed espresse in termini di quantità ad un fluido appena definite.

\subsection{Equazione di continuità ad un fluido}

Ricordiamo che l'equazione di continuità per una singola specie è $\frac{\partial n_\alpha}{\partial t} = -\vec \nabla\cdot (n_\alpha \vec u_\alpha)$. Partiamo da quest'ultima, moltiplicandola per $m_\alpha$ e facciamo una somma sulle specie per ottenere una nuova equazione di continuità a singolo fluido:

\[ \frac{\partial}{\partial t} \left(\sum_\alpha m_\alpha n_\alpha\right) + \vec \nabla \cdot \left(\sum_\alpha m_\alpha n_\alpha \vec u_\alpha \right) = 0 \]

\begin{equation}\label{continuità a un fluido}
	\boxed{\ \frac{\partial\rho}{\partial t} + \vec \nabla \cdot (\rho \vec u) = 0\ }
\end{equation}

In modo del tutto simile potremmo definire l'equazione di continuità di carica:
\begin{equation}\label{continuità di carica a un fluido}
	\boxed{\ \frac{\partial\rho_c}{\partial t} + \vec \nabla \cdot \vec J = 0\ }
\end{equation}


%
%
\vspace{0.5cm}
\subsection{Equazione del moto ad un fluido}

Riprendiamo anche l'eq. del moto per l'impulso trovata per i due fluidi in eq.(\ref{moto}) e la moltiplichiamo per $m_\alpha$, in seguito andremo di nuovo a sommare sulle specie:

\[\frac{\partial(m_\alpha n_\alpha \vec u_\alpha)}{\partial t} + \vec\nabla \cdot (n_\alpha m_\alpha \langle \vec v \vec v\rangle_\alpha) - q_\alpha n_\alpha (\vec E + \frac{1}{c}\vec u_\alpha\times \vec B)\ =\ 0\]

Limitandoci al primo termine notiamo che sommando su $\alpha$ quello che succederebbe banalmente è che in $\partial_t\, \sum_\alpha\, m_\alpha\, n_\alpha\, \vec u_\alpha\, =\, \partial_t\, (\rho \vec u)$ riconosciamo le definizioni che abbiamo dato in precedenza per densità di massa di singolo fluido e velocità media.

Per quanto riguarda il secondo termine ricordiamo che in passato abbiamo usato $\mathbf{P}_\alpha=n_\alpha m_\alpha \langle \vec w \vec w \rangle_\alpha$, con $\vec v = \vec u + \vec w$, ma ora noi vogliamo una pressione di singolo fluido rispetto al centro di massa che andremo a sommare per trovare il tensore di pressione $\mathbf{P} = \sum_\alpha \mathbf{P}_\alpha = p\mathcal{I} + \Pi$:

\[\begin{aligned}
	\mathbf{P}_\alpha\ 
	&=\ n_\alpha m_\alpha \langle \vec w \vec w \rangle_\alpha\ =\\ 
	&=\ n_\alpha\, m_\alpha\ \langle (\vec v - \vec u) (\vec v - \vec u) \rangle_\alpha\ =\ n_\alpha m_\alpha [\langle \vec v\vec v\rangle_\alpha - \langle \vec v\vec u\rangle_\alpha - \langle \vec u\vec v\rangle_\alpha + \langle \vec u\vec u\rangle_\alpha]\ = \\
	&=\ n_\alpha\, m_\alpha\ [\langle \vec v\vec v\rangle_\alpha - \vec u_\alpha\vec u - \vec u\vec u_\alpha + \vec u\vec u]\ \ .
\end{aligned}\]

Usiamo $\langle \vec v\rangle = \vec u_\alpha$ e calcoliamo $\mathbf{P}=\sum_\alpha \mathbf{P}_\alpha$:

\[ \mathbf{P} = \sum_\alpha \mathbf{P}_\alpha = \sum_\alpha n_\alpha m_\alpha \langle \vec v\vec v\rangle_\alpha - \underset{\rho\vec u\vec u}{\underbrace{\left(\sum_\alpha n_\alpha m_\alpha \vec u_\alpha\right)\vec u}} - \underset{\rho\vec u\vec u}{\underbrace{\cancel{\vec u\left(\sum_\alpha n_\alpha m_\alpha \vec u_\alpha\right)}}} + \underset{\rho\vec u\vec u}{\underbrace{\cancel{\sum_\alpha n_\alpha m_\alpha  \vec u\vec u}}} \]

Vogliamo riprendere ora l'eq. del moto ad un fluido moltiplicata per la massa, scritta ad inizio di questo paragrafo, e sommarla sulle due specie.

Ricordando che $\rho \vec u = \sum_\alpha n_\alpha m_\alpha\vec u_\alpha$ vediamo che nell'equazione di inizio paragrafo il primo termine diventa $\partial_t (\rho \vec u)$ mentre il secondo termine può essere ricondotto al tensore di pressione che abbiamo appena espanso, quindi i primi due termini assieme sono:

\[ \sum_\alpha n_\alpha m_\alpha \langle \vec v \vec v \rangle_\alpha\ =\ \mathbf{P}\, +\, \rho \vec u \vec u\]

\[\begin{aligned} \frac{\partial}{\partial t} (\rho\vec u) + \vec \nabla\cdot (\rho\vec u\vec u) + \vec \nabla \cdot \mathbf{P}&=\ \rho\,  \frac{\partial \vec u}{\partial t} + \cancelto{\text{eq. di continuità}}{\vec u  \frac{\partial\rho}{\partial t} + \vec u\vec\nabla\cdot (\rho\vec u)} + \rho(\vec u\cdot \vec \nabla)\vec u + \vec \nabla \cdot \mathbf{P} = \\ &=\ \rho \, \left[ \frac{\partial\vec u}{\partial t} +(\vec u\cdot \vec \nabla)\vec u \right] \,+\, \vec \nabla \cdot \mathbf{P}
	\ =\ \rho\, \frac{D\vec u}{Dt}\, +\, \vec \nabla \cdot \mathbf{P}\end{aligned}\]

Per il terzo termine, quello di forza:

\[\sum_\alpha q_\alpha n_\alpha \left(\vec E + \frac{1}{c}\vec u\times \vec B\right)\ =\ \underset{\rho_c}{\underbrace{\left(\sum_\alpha q_\alpha n_\alpha \right) }} \vec E + \underset{\vec j}{\underbrace{\left( \sum_\alpha q_\alpha n_\alpha \vec u_\alpha \right)}}\times \vec B/c\ =\  \rho_c \vec E + \frac{1}{c}\,\vec j\times \vec B\]

dove $\rho_c$ è uno sbilancio di carica quindi è vicino a zero, ma non nullo. In definitiva, aver sommato le due equazioni del moto fluide ci ha portato alla forma in eq.(\ref{moto MHD}), un'equazione del moto fluida a singolo fluido:

\begin{equation}\label{moto MHD}
	\boxed{\rho\, \frac{D \vec u}{D t} + \vec\nabla\cdot\mathbf{P}-\rho_c \vec E - \frac{1}{c}\,\vec j\times \vec B\ =\ 0}
\end{equation}


%
%
\vspace{0.5cm}
\subsection{Costruiamo la magnetoidrodinamica}

Riportiamo l'equazione del moto fluida di singola specie perché andremo ora a torturarla:

\[\frac{\partial (n_\alpha \vec u_\alpha)}{\partial t} +  \nabla\cdot \left[n_\alpha \left( \bar u _\alpha \bar u_\alpha + \langle \vec w_\alpha \vec w_\alpha \rangle \right) \right] - n_\alpha\frac{q_\alpha}{m_\alpha} \left(\vec E + \frac{\vec u_\alpha \times \vec B}{c}\right)\ =\ 0\ \ .\]

Ora si nota che siamo partiti da 2 equazioni quindi possiamo ricavarne altrettante, c'è ancora dell'informazione da estrarre dall'eq.(\ref{moto}) del moto qui sopra riportata. Andremo infatti a fare sia la somma che la differenza di 


Prendiamo di nuovo la medesima equazione e stavolta la moltiplichiamo per $q_\alpha$, per poi sommare su $\alpha$:

\[\underset{\partial \vec j/\partial t}{\underbrace{\sum_\alpha \frac{\partial(q_\alpha n_\alpha \vec u_\alpha)}{\partial t} }} + \sum_\alpha q_\alpha \nabla \cdot (n_\alpha \langle \vec v \vec v\rangle_\alpha)\, -\, \sum_\alpha\, q^2_\alpha \,\frac{n_\alpha}{m_\alpha} (\vec E + \frac{1}{c}\vec u_\alpha\times \vec B)\ =\ 0
\]

ora questa diversa forma da cui partire ci condurrà all'equazione di Ohm generalizzata. 

Per il secondo termine torniamo a giocare con la definizione di velocità peculiare:

\[ \begin{aligned} \langle \vec v \vec v\rangle\ =\ \langle (\vec u + \vec w)(\vec u + \vec w) \rangle\ &=\ \langle \vec w \vec w\rangle_\alpha\, +\, \langle \vec w \vec u\rangle_\alpha\, +\, \langle \vec u \vec w\rangle_\alpha \, +\, \vec u\vec u\ =\\ &=\ \langle \vec w \vec w\rangle_\alpha \,+\, (\vec u_\alpha - \vec u)\, \vec u \, +\, \vec u\, (\vec u_\alpha - \vec u)\, +\, \vec u\vec u\ =\\
	&=\ \langle \vec w \vec w\rangle_\alpha\, +\, 2\vec u_\alpha\cdot \vec u\, +\, \vec u\cdot \vec u\ \ .
\end{aligned}\]

\subsubsection{Equazione dell'energia di singolo fluido}
Se ora andiamo a richiamare l'equazione dell'energia del capitolo precedente - che riscriviamo qui sotto - possiamo trovare una chiusura (?) $\sum_\alpha \frac{m_\alpha}{2}$:

\[ \frac{\partial \color{blue}{\varepsilon}}{\partial t}\ =\ -\vec \nabla\cdot \color{blue}{\vec Q} \color{black} + \vec E\cdot\color{blue}{\vec J} \]

\[ \frac{\partial}{\partial t} \color{blue}{\int \frac{1}{2} m v^2 f d\vec v} \color{black}\, +\,  \frac{\partial}{\partial\vec x} \color{blue}{\int\frac{1}{2}m v^2 \vec v f d\vec v} \color{black}\, -\, \color{blue}{q\, n\,\vec w} \color{black} \cdot \vec E\ =\ 0\]

che quindi facciamo diventare:

\[\sum_\alpha \frac{m_\alpha}{2} \frac{\partial}{\partial t} (n_\alpha \langle v^2 \rangle_\alpha)\, +\,\frac{1}{2} \sum_\alpha m_\alpha \nabla \cdot (n_\alpha \langle v^2 \vec v\rangle_\alpha)\, -\, \sum_\alpha q_\alpha n_\alpha \vec w_\alpha \cdot \vec E\ =\ 0\]


Per il terzo termine sappiamo ora di poter sostituirvi la definizione della corrente, $\sum_\alpha q_\alpha n_\alpha \vec w_\alpha\, \equiv\, \vec j$. 

Guardando al primo termine vediamo che questo rappresenta ora la variazione dell'energia interna $\ \partial_t\, \left( \rho\, u^2\, /\, 2\ +\ 3\, p\, /\, 2 \right)\ $, sappiamo che $P_\alpha = \frac{1}{3} m_\alpha n_\alpha \langle \vec w\vec w \rangle = n_\alpha T_\alpha$ e per magia ci ritroviamo in mano:

\[ \frac{\partial}{\partial t} \,\left[\frac{1}{2} \sum_\alpha n_\alpha m_\alpha \langle v^2 \rangle_\alpha \right]\ =\ \frac{\partial}{\partial t } \left[\frac{1}{2} \rho u^2 + \frac{3}{2} \sum_\alpha n_\alpha T_\alpha\right]\ \ . \]

Mentre, per il secondo termine vediamo che dato che la velocità della singola specie è presa rispetto alla velocità del singolo fluido, allora 
\[\langle \vec w\rangle_\alpha\,=\, \vec u_\alpha + \vec u \quad \iff \quad \vec v_\alpha = \vec u + \vec w_\alpha\ \ ,\]
quindi non è più a media nulla. Cerchiamo ora di calcolare $\langle v^2 \vec v\rangle_\alpha$:

\[ \begin{aligned}\langle\, v^2\, \vec v\,\rangle_\alpha\ &=\ \langle\ [\ (\vec w +\vec u)\cdot(\vec w + \vec u)\ ]\ (\vec w + \vec u)\ \rangle_\alpha \\ 
	&=\ \langle w^2 \vec w\rangle_\alpha\, +\,\langle \vec w^2 \vec u\rangle_\alpha\, +\, \langle u^2 \vec w\rangle_\alpha\, +\, u^2\vec u\, +\, 2\langle (\vec u\cdot \vec w)\,\vec w\rangle_\alpha\, +\, 2\,\langle (\vec u\cdot \vec w)\,\vec u\rangle_\alpha
\end{aligned} \]

e - a questo punto - cerchiamo un modo intelligente di ridurre questi termini, in particolare riconoscendo i flussi di calore di singola specie $\vec q_\alpha$, le pressioni di singola specie $p_\alpha$ e ricordando le definizioni di $\rho$, $\rho\vec u$ e $\mathbf{P}$ andiamo a vedere cosa diventano questi termini assieme alla somma $\sum_\alpha m_\alpha/2$ ma senza considerare per il momento il gradiente:


\[ (\vec w\cdot \vec u) \, \vec w\ =\ \vec w \vec w\cdot \vec u\ \ \Rightarrow\ \ \sum_\alpha n_\alpha m_\alpha \, \langle \vec w \vec w \rangle_\alpha \cdot \vec u\ =\ \mathbf{P}\cdot \vec u\ \ ,\]

dove $\vec w \vec w=w_iw_j$ è da intendersi come tensore.

\[ \frac{1}{2} \sum_\alpha n_\alpha m_\alpha\, \langle w^2 \vec w\, \rangle_\alpha \ \equiv \sum_\alpha \vec q_\alpha\ =\ Q\]

\[ \frac{1}{2} \sum_\alpha n_\alpha m_\alpha \langle w^2 \vec u \rangle_\alpha\ =\ \frac{1}{2} \sum_\alpha n_\alpha m_\alpha \langle w^2 \rangle _\alpha \vec u\ =\ \frac{3}{2} \sum_\alpha p_\alpha \vec u\ =\ \frac{3}{2} \sum_\alpha n_\alpha T_\alpha \vec u\]

\[ \frac{1}{2} \sum_\alpha n_\alpha m_\alpha u^2\ =\ \frac{1}{2} \rho u^2 \vec u\]

\[ \langle u^2 \vec w_\alpha \rangle = u^2 \langle \vec w \rangle_\alpha\ =\ u^2\ \langle\, \vec u_\alpha - u\, \rangle\ =\ \cancel{\rho\vec u} - \cancel{\rho\vec u}\ =\ 0\ \ . \]

Quindi, dopo tutti questi rimaneggiamenti abbiamo che per il \underline{secondo termine dell'equazione dell'energia}, in definitiva:

\[ \frac{1}{2} \sum_\alpha m_\alpha \nabla\, \cdot\, ( n_\alpha\, \langle v^2 \vec v \, \rangle_\alpha )\ =\ \nabla \cdot \left(\frac{1}{2}\rho u^2 \vec u + \frac{3}{2}\, p\, \vec u\, + \mathbf{P}\cdot \vec u\, +\vec q\,\right) \]


Otteniamo così l'eq dell'\textbf{energia per singolo fluido} (in caso \emph{collision-less}):

\begin{equation}
	\label{Energia singolo fluido}
	\boxed{\frac{\partial}{\partial t} \left(\frac{1}{2}\rho u^2 +\frac{3}{2} p\right) + \vec\nabla \cdot \left( \frac{1}{2}\rho u^2 \vec u + \frac{3}{2}\rho\vec u  + \mathbf{P}\cdot\vec u + \mathbf{q}\right) \ =\ \vec j\cdot \vec E}
\end{equation}


\subsubsection{Plasma di elettroni e protoni}
Torniamo quindi all'equazione del moto e per semplificare i conti assumiamo che gli ioni siano protoni così da avere carica $+e$:

\[\begin{cases}
	\vec j = \sum_\alpha q_\alpha n_\alpha \vec u_\alpha = n_e (\vec u_i-\vec u_e)\, ;\\
	\vec u = \vec u_i + \frac{m_e}{m_i} \vec u_e
\end{cases}
\quad \quad \rightarrow \quad \quad 
\begin{cases}
	\vec u_i = \vec u +\frac{m_e}{m_i} \frac{\vec j}{n_e}\\
	\vec u_e = \vec u - \frac{j}{n_e}
\end{cases}\]

\[ \frac{\partial \vec j}{\partial t} + \sum_\alpha q_\alpha \nabla\cdot (n_\alpha \langle v v\rangle_\alpha ) - \sum_\alpha \frac{n_\alpha}{m_\alpha} q_\alpha^2 (\vec E + \frac{1}{c} \vec u_\alpha \times \vec B)\ =\ 0 \]

Per il secondo termine possiamo 'portare fuori' la divergenza e sviluppiamo:

\[\begin{aligned}
	\sum_\alpha n_\alpha q_\alpha \langle vv\rangle_\alpha &= ne \langle \vec v\vec v\rangle_i - ne \langle \vec v\vec v\rangle_e\ =\\
	&= ne\left[ \langle\, [(\vec v_i - \vec u)+\vec u]\,[(\vec v_i - \vec u)+\vec u]\,\rangle_i - \langle\, [(\vec v_e - \vec u)+\vec u]\, [(\vec v_e - \vec u)+\vec u]\,\rangle_e \right] =\\
	&= n e \left( \underset{\text{press. singola specie } \langle w w \rangle_i}{\underbrace{\langle (\vec v_i -\vec u)\, (\vec v_i - \vec u) \rangle_i}} + \cancelto{0}{\vec u\vec u} + \langle (\vec v_i - \vec u)\, \vec u \rangle_i + \underset{\vec u\langle v_i - \vec u\rangle_i=\vec u(\vec u_i-\vec u)}{\underbrace{\langle \vec u (\vec v_i - \vec u)\rangle_i}}\right) - ne\,(\dots)_e\,
\end{aligned}\]

da cui in definitiva considerando i termini che si elidono fra $+ne(\dots)$ e $-ne(\dots)$ si trova che il secondo termine a meno del gradiente diviene:

\[ \sum_\alpha n_\alpha q_\alpha \langle vv\rangle_\alpha\ =\ \frac{e}{m_i} p_i - \frac{e}{m_e} p_e + (\vec j\,\vec u +\vec u\, \vec j)\ \left(1 + \cancelto{\text{trascurabile in prima analisi}}{\frac{m_e}{m_i}}\right) \]

Quindi nell'equazione del moto abbiamo:

\[ \frac{\partial \vec j}{\partial t} + \nabla \cdot \left[ e\, \left(\frac{p_i}{m_i}- \frac{p_e}{m_e}\right)\, +\,  \vec j \vec u + \vec u\vec j\ \right] \,-\, \frac{n e^2}{m_e}\, \left( \vec E + \frac{\vec u\times \vec B}{c} - \frac{\vec j\times \vec B}{n\, e\, c} \right)\ =\ \eta\dots \text{(termini resistivi)}\ \ . \]

Prendiamo infine il termine finale dell'equazione del moto, quello con $q_\alpha^2$:

\[\begin{aligned}
	\sum_\alpha \frac{n_\alpha}{m_\alpha} q^2_\alpha (\vec E + \frac{\vec u_\alpha\times \vec B}{c})\ &=\ n\, e^2\, \left(\cancel{\frac{1}{m_i}} + \frac{1}{m_i}\right)\vec E\, +\, n\,e^2\, \left(\sum_\alpha \frac{\vec u_\alpha}{m_\alpha}\right)\times \vec B/c\ =  \\ &= \frac{n\, e^2}{m_e} \vec E\, +\, \frac{n^2_e}{m_e}\, \left(\vec u \, + \,\frac{\vec j}{m_e}\right)\times \frac{\vec B}{c}
\end{aligned}\]

dove, per il termine in parentesi - nell'ultimo passaggio - si sono usate le espressioni di $\vec u_i$ e $\vec u_e$ introdotte ad inizio paragrafo per sviluppare:

\[ \sum_\alpha\, \frac{\vec u_\alpha}{m_\alpha}\ =\ \frac{\vec u}{m_i}\, +\, \frac{m_e \, \vec j}{m_i^2\, n_e}\, +\, \frac{\vec u}{m_e}\, -\,\frac{\vec j}{m_e\, n_e}\ =\ \frac{1}{m_e}\, \left[ \vec u\, \left(1+\frac{m_e}{m_i}\right)\, -\, \frac{\vec j}{n_e}\, \left(1+\frac{m_e}{m_i}\right) \right] \]

Otteniamo così per il terzo termine:

\[\sum_\alpha \frac{n_\alpha q^2_\alpha}{m_\alpha} \left(\vec E + \frac{\vec u_\alpha \times \vec B}{c}\right) \approx \frac{n^2_e}{m_e} \left(\vec E + \frac{1}{c} \vec u\times \vec B - \frac{1}{n\, e\, c} \vec j\times \vec B\right)
\]

Infine eccoci \textbf{finalmente} giunti all'\underline{equazione di Ohm generalizzata}:

\begin{equation}
	\label{eq di Ohm generalizzata}
	\boxed{\ \frac{\partial \vec j}{\partial t}\, +\, \vec \nabla\cdot \left[\vec j\,\vec u\, +\, \vec u\, \vec j\, -\, \frac{p_e}{e}\right]\, -\,\frac{n\, e^2}{m_e} \left(\vec E + \frac{1}{c} \vec u\times \vec B - \frac{1}{n\, e\, c} \vec j\times \vec B\right) \ =\ 0\ }
\end{equation}

\vspace{0.5cm}

\subsubsection{Altre considerazioni per costruire MHD}

\paragraph{Faraday in MHD : } Consideriamo la legge di Faraday $\frac{1}{c}\, \frac{\partial \vec B}{\partial t}\, =\, -\vec \nabla \times \vec E$, dimensionalmente abbiamo che $B\, /\, c\tau \sim E\, / L$ quindi anche $\frac{E}{B}\, \sim\, \frac{L}{c\, \tau}\,\sim\, \frac{u}{c}$ così che abbiamo una stima di una velocità caratteristica in funzione di lunghezza e tempo caratteristici $u\sim L/\tau$.

Nella MHD ideale andiamo a supporre che $\vec E\, +\, \frac{\vec u\times \vec B}{c}\, =\, 0$ quindi, in unità $c=1$, vediamo che l'equazione di Faraday-Maxwell in MHD ideale assume la forma:

\begin{equation}
	\label{Faraday MHD ideale}
	\boxed{\frac{\partial \vec B}{\partial t}\, =\, \vec \nabla \times \left(\vec u\times \vec B\right)}\ \ \ .
\end{equation}

\paragraph{Ampére in MHD : } Dalla legge di Ampère sappiamo che $\frac{1}{c}\frac{\partial E}{\partial t}\, =\, \vec \nabla \times \vec B\, -\, \frac{4\pi}{c} \vec j$, andiamo a fare ancora un'analisi dimensionale per capire quali termini dominino:

\[ \Rightarrow\quad \frac{|\partial_t\, \vec E|}{c\,|\vec \nabla\, \times\, \vec B|}\ \sim\ \frac{E\, L}{c\, \tau\, B}\ \sim\ \left( \frac{u}{c}\right)^2 \]

questa ci permette di dire che nel caso in cui $u\ll c$, allora il termine $\partial_t \vec E$ è trascurabile rispetto al termine $\vec \nabla \times\vec B$, quindi l'equazione di \underline{Ampère-Maxwell in MHD} diventa: 
\begin{equation}
	\label{Ampére MHD}
	\boxed{\vec \nabla\times \vec B\ \approx\ \frac{4\pi}{c}\vec j}\ \ ,
\end{equation}
quindi, dimensionalmente, diciamo anche che $B/L\, \sim\, j / c$ ed osserviamo che la corrente $\vec j$ è associata al campo e non legata ad $\vec u_e$ e $\vec u_i$.

\paragraph{Equazione del moto in MHD : } Stiamo pian piano costruendo questo nostro modello della magnetoidrodinamica, riconsideriamo ora l'eq. del moto eq.(\ref{moto MHD}), per la quale eravamo rimasti con la forma $\ \rho\,\frac{D\, \vec u}{D\, t}\,=\,- \nabla p + \rho_c\, \vec E + \vec j\times \vec B\, /\, c$ . In questo caso non possiamo dire che il termine $\rho_c\, \vec E$ si annulli solo per la quasi-neutralità, $n_e\, \simeq\, n_i$, ma possiamo considerarlo come un'approssimazione del termine $\frac{1}{c} \vec j\times \vec B$ , infatti:

\[ \frac{\rho_c\, |\vec E|}{|\vec j\times \vec B|\, /\, c}\ \sim \ \frac{E^2}{L}\,\frac{c}{j\,B}\ \sim\ \frac{E^2}{B^2}\ \sim\ \left(\frac{u}{c}\right)^2\ \ll\ 1 \ \ .\]

quindi nell'\underline{equazione del moto} possiamo trascurare $\rho_c\, \vec E$ e scrivere la versione \underline{semplificata}:

\begin{equation}
	\label{moto MHD semplificato}
	\boxed{\rho\, \frac{D\, \vec u}{D\, t}\, +\, \vec\nabla\cdot \mathbf{P} \, -\, \frac{1}{c}\vec j\times\vec B\ =\ 0}\ \ .
\end{equation}

Nel caso stazionario $\vec \nabla \cdot \mathbf{P}\, =\, \vec j\times \vec B\, /\, c$ , da cui dimensionalmente: $P\,/\,L\,\sim\, j\,B\,/\,c\,\sim\, B^2\,/\,L$ .

\paragraph{Altre definizioni importanti : } Ancora sulla linea del fornire altre rifiniture per completare il nostro modello MHD, definiamo un'altra velocità caratteristica del plasma diversa dalla velocità del suono, cioè la \underline{velocità di Alfvén}, che partendo da $\rho \frac{\partial u}{\partial t}\sim \frac{j}{c}\, B= \frac{1}{4\pi}\frac{B^2}{L}$ e scrivendo $\rho \frac{\partial u}{\partial t}\sim \rho \frac{u_a}{\tau}$, si definisce come:

\begin{equation}
	\label{velocità di Alfvén}
	\boxed{\,v_a^2\ =\ \frac{B^2}{4\pi \rho}\,}
\end{equation}

Questa è la velocità con cui si propagano le onde di Alfvén ma anche pi
ù in generale la velocità massima di propagazione delle perturbazioni magnetoidrodinamiche.

Altra quantità che solitamente si definisce in ambito di fusione soprattutto è il \underline{parametro $\beta$} del plasma, definito come:

\begin{equation}
	\label{beta parameter}
	\boxed{\,\beta\ \equiv\ \frac{8\, \pi\, p}{B^2}\,}
\end{equation}

\paragraph{Equazione di Ohm generalizzata : } Ripartiamo dall'equazione di Ohm generalizzata, ma riarrangiata rispetto ad eq.(\ref{eq di Ohm generalizzata}) e questa volta non ideale quindi inserendo il termine resistivo:

\begin{equation}
	\label{Ohm come dio comanda}
	\underset{\text{termini in situazione ideale}}{\underbrace{\vec E + \frac{\vec u\times \vec B}{c}}}\, -\, \underset{\text{termine resistivo}}{\underbrace{\eta\vec j}}\ =\ \underset{\text{termini di scala ionica}}{\underbrace{\frac{1}{n e c} \,\vec j\times \vec B\, -\, \frac{1}{n\, e} \nabla\cdot \mathbf{P}_e}}\, +\, \underset{\text{termini di scala elettronica}}{\underbrace{\frac{m_e}{n\, e^2}\, \left[\frac{\partial \vec j}{\partial t}\, +\, \nabla \cdot (\vec j\otimes \vec u + \vec u \otimes \vec j)\right]}}\ \ .
\end{equation} 

Ci fermiamo ad analizzare la rilevanza dei termini, innanzitutto supponendo che $\vec \nabla\cdot \mathbf{P}_e \sim \vec \nabla p_e\ \iff \ |\vec \nabla\cdot \mathbf{\Pi}|\ll\vec \nabla p_e$ , cioè andiamo a trascurare l'anisotropia ed andiamo a studiare la rilevanza del termine rispetto all'altro termine di scala ionica:

\[ \frac{\ |\vec\nabla p_i|\,/\, n\,e\ }{\ |\vec u\times \vec B|\,/\,c\ }\ \sim\ \frac{m_i\, n\, v_i^2}{L\, n\, e}\frac{c}{u\, B}\ \underset{v_i\sim u}{\sim}\ \frac{m_i\, c\, u}{e\, B\, L}\ \sim \ \frac{\omega}{\omega_{cicl.}}\ \ . \]

La stima fatta ci porta ad avere una condizione per semplificare l'eq. di Ohm in funzione della frequenza di ciclotrone degli ioni $\omega_{cicl.}\,=\,\frac{e\, B}{m_i\, c}$, ma ci piacerebbe riportare l'analisi in termini di frequenza di plasma $\omega_{p,\, i}\, \sqrt{4\pi\, n\, e^2\, /\, m_i}$ , è qui che torna utile la velocità di Alfvén: $v_A\,\equiv\, B\, /\, \sqrt{4\pi\, n\, m_i}$ .
Osserviamo che:
\[\frac{e\, B}{m_i}\ =\ \frac{B}{\sqrt{4\pi n\, m_i}} \sqrt{\frac{4\pi \, n\, e^2}{m_i}}\ =\ v_A\, \omega_{p,\, i}\ \ ,\]
quindi per la frequenza di ciclotrone vale che $\omega_{cicl.}\, =\, \frac{v_A\, \omega_{p,\, i}}{c}\, =\, \frac{v_A}{\delta_i}$, dove $\delta_i\, =\, c\, /\, \omega_{p,\, i}$ è la lunghezza pelle degli ioni che avevamo già definito in passato come $\delta_s = \displaystyle\sqrt{\frac{m_s}{4\pi\, n\, q_s^2}}\, c$.

Cosa possiamo finalmente dire sulla predominanza dei termini? che nell'ipotesi di $v_A\sim u$, tipica in magnetoidrodinamica, il plasma trasporta energia alla stessa velocità con cui si propagano i segnali magnetici:

\[\Rightarrow \quad \frac{\omega}{\omega_{cicl.}}\ =\ \frac{\omega \delta_i}{v_A}\, \sim\, \frac{v_A}{u}\, \frac{\delta_i}{L}\ \ .\]

quindi, finché $L\gg \delta_i$, ovvero finché non ci troviamo su scale ioniche, il termine $\vec \nabla p_e$ è solo una correzione del termine $\vec u\times \vec B$.

Analogamente, si può mostrare che negli ultimi 2 termini dell'equazione di Ohm in forma eq.(\ref{Ohm come dio comanda}) diventino rilevanti solo a scale elettroniche:

\[ \frac{m_e}{n\, e}\, \left| \frac{\partial \vec j}{\partial t} \right|\, \frac{c}{|\vec u\times \vec B|}\ \sim\ \frac{m_e}{n\, e^2}\, \frac{j}{\tau}\, \frac{c}{u\, B} \]
\[ \underset{\text{Faraday, } B\sim L\,j/\,c}{\Rightarrow} \quad \quad \quad \frac{m_e}{n\, e^2}\, \frac{c^2}{L\, \tau\, u}\ \sim\ \left(\frac{m_e\, c^2}{n\,e^2}\right)\, \frac{1}{L^2}\ \sim\ \frac{\delta_e^2}{L^2}\ \ .\]

In definitiva quel che capiamo è che la versione più semplificata possibile risulti essere:

\begin{equation}
	\label{Ohm scema}
	\boxed{\vec E\, +\, \frac{\vec u\times \vec B}{c}\ =\ 0}\ \ ,
\end{equation}

che è la nostra \underline{equazione di Ohm in MHD ideale}, se teniamo il termine resistivo abbiamo:

\begin{equation}
	\label{Ohm resistiva}
	\boxed{\vec E\, +\, \frac{\vec u\times \vec B}{c}\ =\ \eta \vec j}\ \ .
\end{equation}

detta \underline{equazione di Ohm resistiva}.

Un'altra versione accennata a lezione è la \emph{'electron MHD Ohm's law'} $\vec E\, +\, \frac{\vec u_e\times \vec B}{c}\ =\ 0$ che però non saprei in che condizioni utilizzare.


%
%
\subsubsection{Ancora sull'equazione dell'energia}

Torniamo sull'equazione dell'energia per singolo fluido, in eq.(\ref{Energia singolo fluido}), in modo da capire in che condizioni ridurci ad una politropica:

\[ \frac{\partial}{\partial\, t} \left(\frac{1}{2}\,\rho\, u^2\, +\, \frac{3}{2}\, p\right)\ + \ \vec\nabla\cdot\left[\frac{1}{2}\,\rho\, u^2\,\vec u\, +\, \frac{3}{2}\, p\, \vec u\, +\, \mathbf{P}\cdot\vec u\, +\, \vec q \right]  \ =\ \vec j\cdot \vec E \]

\[ \mathbf{P}\cdot \vec u\ =\ \left( p\mathbf{I} + \mathbf{\Pi} \right)\cdot \vec u \ \simeq\ p\mathbf{I}\cdot \vec u \]

\begin{equation}
	\label{Energia singolo fluido v2}
	\frac{\partial}{\partial\, t} \left(\frac{1}{2}\,\rho\, u^2\, +\, \frac{3}{2}\, p\right)\ + \ \vec\nabla\cdot\left[  \underset{u^2\,(\rho\vec u)}{\underbrace{\frac{1}{2}\,\rho\, u^2\,\vec u\,}} +\, \frac{5}{2}\, p\, \vec u\, +\, \vec q \right]  \ =\ \vec j\cdot \vec E
\end{equation}

\[ \frac{1}{2}\, u^2\,\cancel{\frac{\partial \rho}{\partial\, t}}\, +\, \frac{1}{2}\, \rho\,\frac{\partial (u^2)}{\partial\, t}\, +\,  \frac{3}{2}\, \frac{\partial p}{\partial\, t}\ + \ 
\frac{1}{2}\, u^2\, \cancel{\vec\nabla\cdot(\rho\vec u)}\, + \, \frac{1}{2} \left(\rho\vec u\cdot \vec \nabla\right)\, u^2\, +\, \frac{5}{2}\, p\, \vec\nabla\cdot \vec u\, +\, \frac{5}{2} (\vec u\cdot \nabla)\, p \, +\, \vec\nabla\cdot\vec q \ =\ \vec j\cdot \vec E \]

Nell'ultimo passaggio abbiamo semplicemente sviluppato questa nostra forma approssimata raggiunta in eq.(\ref{Energia singolo fluido v2}).

Vogliamo ora sottrarvi l'equazione (di seguito) del moto, in eq.(\ref{moto MHD}) moltiplicata per $\vec u\cdot$:

\[ \vec u\cdot \rho\, \frac{\partial \vec u}{\partial t}\, +\, \rho\vec u\cdot \nabla \left(\frac{u^2}{2}\right) \,+\, \vec u\cdot\vec\nabla p\ =\ \frac{1}{c}\,\vec u\cdot \left(\vec j\times \vec B \right) \]

Così da ottenere:

\[ \frac{3}{2}\, \frac{\partial p}{\partial t}\, +\, \frac{5}{2}\, p\, \vec \nabla\cdot\vec u\, +\, \frac{3}{2}\left(\vec u\cdot \vec \nabla\right)\, p\ =\ -\vec\nabla\cdot\vec q \, +\, \vec j\cdot \vec E\, -\, \frac{1}{c}\,\vec u\cdot \left(\vec j\times \vec B \right)\]

Giunti a questo punto possiamo chiederci quali siano le approssimazioni adeguate ad esempio trascurare il flusso di calore (legato al termine $\vec\nabla\cdot\vec q$) o le dissipazioni (effetto Joule, legato al termine $\vec j\cdot \vec E$).

Per invarianza del prodotto misto sotto permutazioni cicliche otteniamo $\vec j\cdot(\vec j\times \vec B)=-\vec u\cdot(\vec B \times \vec j)=-\vec j\cdot (\vec u\times \vec B)$ che vuol dire che possiamo prendere i due termini $+\, \vec j\cdot \vec E\, -\, \frac{1}{c}\,\vec u\cdot \left(\vec j\times \vec B \right)$ e permutare il prodotto per ottenere $\vec j\cdot\vec E'$, per un nuovo $\vec E'=\vec E + \frac{1}{c}\vec u\times \vec B$.

\[\Rightarrow\quad \quad \underset{ \frac{3}{2} \left[\frac{\partial}{\partial t}+\vec u\cdot \vec \nabla\right]\, p }{\underbrace{\frac{3}{2}\, \frac{\partial}{\partial t} p\, +\, \frac{3}{2}\, \left(\vec u\cdot \vec \nabla\right)\, p}}\, +\, \frac{5}{2}\, p\vec \nabla\cdot \vec u\ =\ 0\ \ .\]

\'E proprio in queste condizioni (trascurando flusso di calore e dissipazioni) che troviamo la giustificazione del perché la politropica sia una buona equazione per chiudere il sistema, la \underline{politropica} è $\boxed{p\rho^{-\gamma}=\text{cost.}}$, quindi derivandola:

\[ 0\ \overset{!}{=}\ \frac{D\, \left(p\rho^{-\gamma}\right)}{D\, t}\ =\ \frac{D\, p}{D\, t}\, \rho^{-\gamma} \ +\ p\gamma \rho^{-\gamma -1} \frac{D\, \rho^{-\gamma -1}}{D\, t}\ =\ \underset{=0\ \text{per $\gamma=5/3$ è come la forma trovata sopra}}{\rho^{-\gamma}\, \underbrace{\left[\frac{Dp}{Dt}+\gamma p\vec\nabla\cdot\vec u\right]}}\]






%
%
\subsection{Riassumendo il framework di MHD ideale}
Abbiamo infine un sistema di equazioni per l'MHD ideale che comprende:

\begin{equation}
	\label{sistema MHD}
	\begin{aligned}
		\text{continuità:} \qquad 
		& \frac{\partial \rho}{\partial t} + \vec\nabla \cdot (\rho \vec u) = 0 \\		
		\text{eq.\ del moto:} \qquad 
		& \rho\!\left( \frac{\partial \vec u}{\partial t} + (\vec u \cdot \vec\nabla)\vec u \right)
		= -\vec\nabla p + \frac{1}{4\pi}(\vec\nabla \times \vec B)\times \vec B \\
		\text{Faraday-Maxwell:} \qquad 
		& \frac{\partial \vec B}{\partial t}
		= \vec\nabla \times (\vec u \times \vec B) \\
		\text{politropica:} \qquad 
		& p\,\rho^{-\gamma} = \text{cost.}\\
		\text{campo elettrico:} \qquad 
		& \vec E\ =\ -\frac{1}{c}\vec u\times \vec B\\
		\text{corrente:} \qquad 
		& \vec j\ =\ \frac{c}{4\pi} \, \vec \nabla\times \vec B
	\end{aligned}
\end{equation}

